\section{Statement}
\label{sec:statement}

The primary goal of this research is to advance the state of the art in cloud-native observability by transitioning from ad-hoc tool chaining to a formalized architectural model. To this end, the core contributions of this thesis are defined as follows:

\begin{itemize}
    \item \textbf{The PANOPTES Reference Architecture:} A formalized, vendor-agnostic architectural blueprint for Kubernetes observability. It structurally maps how metrics, logs, and traces must be generated, collected, and semantically correlated to achieve system observability in accordance with control theory principles.
    
    \item \textbf{An Open-Source Implementation Artifact:} A concrete instantiation of the PANOPTES architecture using Infrastructure as Code (Vagrant, Ansible, and Helm). This artifact operationalizes tools such as OpenTelemetry, Prometheus, Loki, and Tempo in a reproducible, deterministic environment, providing a ready-to-use blueprint for practitioners and researchers.
    
    \item \textbf{Empirical Evidence via Chaos Engineering:} A quantitative evaluation of the proposed architecture. Through controlled experiments involving fault injection, this study provides empirical data on the actual resource overhead introduced by the observability stack and its measurable impact on the Mean Time to Detect (MTTD) anomalies, establishing a scientific baseline for future comparisons.

    \item \textbf{The Observability Maturity Model (OMM):} A conceptual framework derived from a Rapid Review of the state-of-the-art, categorizing the evolutionary stages of telemetry management in cloud-native environments. This model serves as a roadmap for organizations transitioning from fragmented monitoring to intelligent, predictive observability.
\end{itemize}